\section{Discusión}

La dimensión que se utilice para medir las diferencias territoriales cambia
por completo la lectura de la convergencia regional en el Perú. El PBI real
per cápita relativo exhibe una relación negativa y estadísticamente
significativa entre el nivel inicial y el cambio acumulado; el IDH, en
cambio, exhibe el signo contrario. Durante 2007--2019, los departamentos se
acercaron entre sí en términos económicos sin que esa aproximación tuviera
una contraparte equivalente en desarrollo humano.
 
Ingreso y bienestar multidimensional rara vez avanzan al mismo compás. Un
territorio puede incrementar su producción por habitante sin que ese
crecimiento se traduzca de inmediato en mejores logros educativos, mejores
condiciones de salud o mayores oportunidades; esa traducción depende de la
provisión de servicios públicos, la infraestructura social, la calidad
institucional y la manera en que cada departamento distribuye internamente
sus recursos. El enfoque de capacidades sitúa precisamente ahí el punto
central: el crecimiento económico es un medio para ampliar las libertades de
la población, no una medida suficiente del desarrollo en sí misma
\parencite{sen1999}.
 
Que el coeficiente de variación del IDH haya disminuido podría interpretarse,
a primera vista, como una aproximación relativa entre departamentos. El signo
positivo del coeficiente beta obliga a matizar esa lectura: los territorios
que partían con niveles más altos también fueron los que más avanzaron. No
hay contradicción entre ambos resultados, porque responden a preguntas
distintas. La convergencia sigma describe cómo evoluciona la dispersión total
de la distribución; la convergencia beta pregunta si la posición inicial
predice el cambio posterior. Una distribución puede volverse menos dispersa
en términos relativos y, al mismo tiempo, conservar una ventaja dinámica para
quienes ya estaban mejor posicionados.
 
Esta distinción recupera un argumento planteado tempranamente por
\textcite{quah1993}: reducir la dinámica regional a un único coeficiente
puede ocultar más de lo que revela. Las densidades kernel, los mapas y las
regresiones no paramétricas de este estudio muestran precisamente eso: el
promedio nacional no resume el proceso departamental. El IDH mejoró en
términos agregados, pero las trayectorias individuales se mantuvieron
desiguales, con una relación positiva entre el punto de partida y el avance
posterior.
 
El coeficiente OLS de $-0.3282$ estimado para el PBI real per cápita
relativo, junto con la velocidad anual implícita de 3.3\%, constituye
evidencia de convergencia económica, aunque de un tipo particular: una
reducción de las posiciones relativas, no una igualación de niveles
absolutos. Los departamentos que partían rezagados mejoraron en promedio con
mayor rapidez que los que partían adelantados, sin que ello implicara el
cierre completo de las brechas productivas hacia el final del periodo.
 
La caída de la desviación estándar del indicador económico —de 0.6193 a
0.4750— respalda esta lectura y confirma una menor dispersión departamental.
Los valores extremos, sin embargo, siguen distantes entre sí, lo que revela
diferencias todavía relevantes en productividad, estructura sectorial y
capacidad de generación de ingresos. La convergencia observada es, en ese
sentido, un proceso parcial más que la desaparición de las desigualdades
regionales.
 
La literatura peruana sobre el tema anticipa buena parte de este patrón: las
conclusiones sobre convergencia dependen del periodo analizado, la variable
elegida y el método de estimación. \textcite{gonzales_trelles2004} ya habían
subrayado el peso de las diferencias territoriales y de la geografía
económica en la evolución regional; \textcite{palomino_rodriguez2019}
incorporaron la dependencia espacial al estudio del crecimiento
departamental; \textcite{pazparedes2023}, por su parte, identificó
trayectorias agrupadas y factores estructurales detrás de los clubes
regionales.
 
Ningún patrón espacial global resultó significativo para los niveles
iniciales ni para los cambios acumulados, lo que indica que los
departamentos con valores parecidos no se concentran de manera sistemática
en el territorio bajo las matrices empleadas. Eso no equivale a decir que el
espacio carezca de importancia: las pruebas LM, los modelos espaciales y la
GWR muestran que su relevancia cambia según el indicador y la definición de
vecindad que se utilice.
 
Para el IDH, los modelos espaciales conservan el signo positivo del
coeficiente inicial sin que los parámetros de error espacial ofrezcan
evidencia suficiente para desplazar al modelo OLS. El ancho de banda de la
GWR se aproxima además a la distancia máxima entre departamentos, y el rango
de los coeficientes locales es mínimo —señales que, en conjunto, apuntan a
una relación predominantemente global entre el nivel inicial del IDH y su
variación posterior.
 
Conviene, no obstante, ser cautelosos con esa lectura. Un ancho de banda tan
elevado puede reflejar una relación genuinamente homogénea en el territorio,
pero también puede ser síntoma de la limitada capacidad de una muestra de 24
departamentos para detectar variaciones locales. El resultado no demuestra,
entonces, que las dinámicas del desarrollo humano sean idénticas en todo el
país; solamente indica que la GWR empleada aquí no logra identificar
diferencias relevantes en la pendiente estimada.
 
La heterogeneidad espacial resulta bastante más visible en el PBI per
cápita relativo. Los coeficientes locales oscilan entre $-0.4857$ y
$-0.2771$ y conservan signo negativo en los 24 departamentos: hay
convergencia económica en todo el territorio, aunque su intensidad varía de
un lugar a otro. Este hallazgo coincide con lo planteado por
\textcite{duran2024}, para quienes los coeficientes globales pueden esconder
diferencias espaciales relevantes en la velocidad y la magnitud de la
convergencia.
 
También merece atención la sensibilidad de los modelos económicos frente a
la matriz espacial utilizada. La banda de distancia inversa no mejora el
criterio de información respecto del OLS, mientras que el SEM basado en
contigüidad queen sí presenta un AIC menor. Algunas pruebas robustas
detectan además dependencia en el componente de error, lo que sugiere que
las perturbaciones económicas comparten factores territoriales omitidos
entre departamentos vecinos —aunque la evidencia no es completamente
estable entre especificaciones.
 
Estas diferencias entre matrices no son un simple problema técnico. La
contigüidad captura vínculos fronterizos directos; la distancia inversa
asume que la interacción decae de forma continua conforme aumenta la
separación geográfica. En un territorio marcado por barreras naturales,
infraestructura desigual y conexiones económicas que no siempre coinciden
con la cercanía física, cada matriz termina representando un mecanismo de
interacción distinto.
 
En términos de política pública, una estrategia centrada únicamente en el
crecimiento económico regional difícilmente cerraría las brechas de
desarrollo humano. La convergencia del ingreso relativo necesita
complementarse con intervenciones en educación, salud, conectividad, empleo
de calidad y acceso a servicios básicos; la asignación de recursos debería
considerar tanto las condiciones iniciales de cada departamento como su
capacidad para transformar el crecimiento productivo en bienestar efectivo.
 
La heterogeneidad detectada en la convergencia económica pone en cuestión,
además, la aplicación uniforme de políticas regionales. Departamentos con
una convergencia más débil podrían requerir instrumentos específicos de
diversificación productiva, infraestructura logística y fortalecimiento del
capital humano, mientras que los vínculos entre territorios vecinos
justifican mecanismos de coordinación interdepartamental en transporte,
mercados, corredores productivos y provisión de servicios.
 
\subsection{Limitaciones del estudio}
 
El tamaño de la muestra es la limitación más evidente: 24 departamentos
restringen la potencia de las pruebas espaciales y de las regresiones
geográficamente ponderadas, por lo que los resultados locales deben leerse
como evidencia exploratoria y no como parámetros estructurales definitivos.
 
El nivel de agregación departamental plantea una segunda limitación, porque
puede esconder desigualdades internas entre provincias, distritos, zonas
urbanas y rurales. Dos departamentos con un promedio similar podrían
albergar distribuciones internas muy distintas; trabajos futuros con
información comparable podrían descender a unidades territoriales más
pequeñas.
 
Trabajar con dos momentos en el tiempo —2007 y 2019— permite estimar una
trayectoria acumulada, pero deja fuera las fluctuaciones intermedias. Un
panel anual habría permitido estimar efectos dinámicos, controlar
características no observadas y examinar cambios en la velocidad de
convergencia a lo largo del periodo.
 
También corresponde señalar una adaptación metodológica puntual: el IDH se
empleó directamente en su escala de cero a uno, en lugar de aplicar la
transformación utilizada en el estudio de referencia. La decisión facilita
la interpretación de los resultados, aunque limita la comparabilidad exacta
con otras réplicas del mismo ejercicio.
 
Por último, los modelos de convergencia incorporan únicamente el nivel
inicial como variable explicativa. Inversión pública, educación,
infraestructura, estructura productiva, urbanización, minería,
transferencias fiscales y calidad institucional podrían explicar una parte
adicional de las diferencias departamentales; su ausencia implica que los
resultados obtenidos identifican asociaciones de convergencia y no
relaciones causales.
